%% kapa: exam=1 round=23 year=2012 period=1 form=A subject=civil_law subject_ko=민법 seq=13 number=13
%% source: Q-Net 감정평가사 기출문제 (cst00309) artlSeq=5052578 file=raw/1008/5052578/2012년 제23회 감정평가사 제1차시험 기출문제(1교시 A형).pdf
%% conversion: kordoc 4.24.51 -> tools/gen1.py (auto); review=proofread-sonnet; flags=none
\begin{question}{13}
만 15세의 미성년자 甲은 친권자 丙의 동의를 얻어 자신의 소유인 X 토지를 乙에게 매도하는 경우, 다음 설명 중 옳은 것을 모두 고른 것은? (다툼이 있으면 판례에 의함)
\begin{framedtext}
ㄱ. 丙의 동의는 명시적으로 뿐만 아니라 묵시적으로도 가능하다. \par
ㄴ. 丙이 동의를 할 때에는 친족회의 동의를 필요로 한다. \par
ㄷ. 丙의 동의는 甲의 법률행위 이전에 하여야 하고 이후에 하는 동의는 효력이 없다. \par
ㄹ. 甲이 丙의 동의가 있는 것처럼 보이기 위하여 소극적인 사술을 쓴 것만으로도 취소권이 배제된다. \par
ㅁ. X 토지가 甲 명의의 매매계약서에 근거하여 乙에게 소유권이전등기가 경료되었다면, 丙의 동의 여부와 관계없이 적법한 등기로 추정된다.
\end{framedtext}
\begin{choices}
  \item ㄱ, ㄹ
  \item ㄱ, ㅁ
  \item ㄴ, ㄷ
  \item ㄴ, ㄹ
  \item ㄷ, ㅁ
\end{choices}
\end{question}
\answer{2}
